\documentclass{article}
\usepackage{graphicx}

\title{A Sample Paper for Testing Integrity Checks}
\author{Test Author}

\begin{document}
\maketitle

\begin{abstract}
This is a test paper that cites both real and fake references.
Our method achieves 95.2\% accuracy on the benchmark.
\end{abstract}

\section{Introduction}
\label{sec:intro}

Deep learning has revolutionized NLP \cite{real_entry_attention, real_entry_bert}.
Recent work by \cite{fake_entry_1} proposed a novel framework.
Cross-lingual approaches \cite{fake_entry_2} have shown promise.

We also reference an undefined key \cite{nonexistent_key}.

\section{Method}
\label{sec:method}

See Figure~\ref{fig:architecture} for an overview.
As shown in Table~\ref{tab:results}, our results are significant.

\input{tables/results}

\includegraphics[width=0.8\textwidth]{figures/architecture.png}

\section{Results}
\label{sec:results}

Our method outperforms baselines as shown in Section~\ref{sec:method}.
We reference a non-existent figure: Figure~\ref{fig:missing}.

\section{Conclusion}
\label{sec:conclusion}

We presented our approach and demonstrated strong results.

\bibliographystyle{plain}
\bibliography{refs}

\end{document}
